\begin{center}
\textbf{БИОИНФОРМАТИКА}

\vspace{0.5cm}

\textit{УДК 57.087.1}

\vspace{1cm}

\textbf{ПРЕДСКАЗАНИЕ МИКРОРНК, ВОВЛЕЧЕННЫХ В РЕПАРПЦИЮ МИОКАРДА, ЧЕРЕЗ СОПОСТАВЛЕНИЕ СЕТЕЙ ГЕН-ГЕННЫХ ВЗАИМОДЕЙСТВИЙ"}

\vspace{1cm}

\textbf{Г. Ж. Осьмак\textsuperscript{a,b,*}, М. С. Козин\textsuperscript{a,b}, О. Г. Кулакова\textsuperscript{a,b}}

\vspace{0.5cm}

\footnotesize
\textsuperscript{a}\textit{Национальный медицинский исследовательский центр кардиологии им. ак. Е.И. Чазова Министерства здравоохранения Российской Федерации, Москва, 121552}

\textsuperscript{b}\textit{Российский национальный исследовательский медицинский университет}

\textit{им. Н.И. Пирогова Министерства здравоохранения Российской Федерации, Москва, 117997}

\vspace{0.5cm}

\textsuperscript{*}\textit{e-mail: german.osmak@gmail.com}

\normalsize
\end{center}

\vspace{0.5cm}

\begin{center}
\textit{Поступила в редакцию}

\textit{После доработки}

\textit{Принята в печать}
\end{center}

\vspace{0.5cm}

\textbf{Аннотация} 

Сердечно-сосудистые заболевания остаются одной из основных причин смертности, во многом из-за ограниченной репарационных возможностей сердечной ткани. В данной работе предложен новый численный подход для выявления микроРНК, потенциально участвующих в процессах репарации миокарда, на основе сопоставления топологических характеристик сетей сигнальных путей и генов-мишеней. В исследовании проанализированы ключевые пути, связанные с восстановлением сердечной и выделены микроРНК, которые могут участвовать в их регуляции. Для сигнального пути HIF-1a: miR651-5p, miR-147a; для VEGF: miR-374b-5p и miR-369-3p; для FGF2: miR-3652, miR-4430 и miR-6134; для TNF-a: miR-23a-3p; для  IGF-1, IL2, PDGF: miR-2117, miR-196b-3p и miR-155-3p, соответственно. Таким образом, предложенный метод позволяет эффективно отбирать перспективные для дальнейшего изучения микроРНК, что открывает новые возможности для разработки терапевтических стратегий, направленных на восстановление функции миокарда.

\vspace{1cm}

\textbf{Ключевые слова:} микроРНК, сигнальные пути, репарация миокарда, сетевой анализ, задача ''О назначениях'', Венгерский алгоритм.

\vspace{1cm}

\textit{Принятые сокращения:} LCC - наибольшая компонента связности